\documentclass[10pt,a4paper]{amsart}
\usepackage{amsmath,amssymb,amsthm,graphicx}
% Retained minimal preamble for the Stacks Project chapter flat.tex (More on Flatness).
% Verbatim excerpt of the chapter preamble: Xy-pic setup (native preamble lines 14-18) and the native
% theorem declarations with the mathematical macros (native preamble lines 154-217). Copied mechanically;
% nothing added, removed or redefined. The wrapper class and the amsmath/amssymb/amsthm/graphicx base
% packages are supplied by the assembler, so the class/externaldocument part of the native preamble is omitted.
% For commutative diagrams we use Xy-pic
\usepackage[all]{xy}

% We use 2cell for 2-commutative diagrams.
\xyoption{2cell}
% Theorem environments.
%
\theoremstyle{plain}
\newtheorem{theorem}[subsection]{Theorem}
\newtheorem{proposition}[subsection]{Proposition}
\newtheorem{lemma}[subsection]{Lemma}

\theoremstyle{definition}
\newtheorem{definition}[subsection]{Definition}
\newtheorem{example}[subsection]{Example}
\newtheorem{exercise}[subsection]{Exercise}
\newtheorem{situation}[subsection]{Situation}

\theoremstyle{remark}
\newtheorem{remark}[subsection]{Remark}
\newtheorem{remarks}[subsection]{Remarks}

\numberwithin{equation}{subsection}

% Macros
%
\def\lim{\mathop{\mathrm{lim}}\nolimits}
\def\colim{\mathop{\mathrm{colim}}\nolimits}
\def\Spec{\mathop{\mathrm{Spec}}}
\def\Hom{\mathop{\mathrm{Hom}}\nolimits}
\def\Ext{\mathop{\mathrm{Ext}}\nolimits}
\def\SheafHom{\mathop{\mathcal{H}\!\mathit{om}}\nolimits}
\def\SheafExt{\mathop{\mathcal{E}\!\mathit{xt}}\nolimits}
\def\Sch{\mathit{Sch}}
\def\Mor{\mathop{\mathrm{Mor}}\nolimits}
\def\Ob{\mathop{\mathrm{Ob}}\nolimits}
\def\Sh{\mathop{\mathit{Sh}}\nolimits}
\def\NL{\mathop{N\!L}\nolimits}
\def\CH{\mathop{\mathrm{CH}}\nolimits}
\def\proetale{{pro\text{-}\acute{e}tale}}
\def\etale{{\acute{e}tale}}
\def\QCoh{\mathit{QCoh}}
\def\Ker{\mathop{\mathrm{Ker}}}
\def\Im{\mathop{\mathrm{Im}}}
\def\Coker{\mathop{\mathrm{Coker}}}
\def\Coim{\mathop{\mathrm{Coim}}}

% Boxtimes
%
\DeclareMathSymbol{\boxtimes}{\mathbin}{AMSa}{"02}

%
% Macros for moduli stacks/spaces
%
\def\QCohstack{\mathcal{QC}\!\mathit{oh}}
\def\Cohstack{\mathcal{C}\!\mathit{oh}}
\def\Spacesstack{\mathcal{S}\!\mathit{paces}}
\def\Quotfunctor{\mathrm{Quot}}
\def\Hilbfunctor{\mathrm{Hilb}}
\def\Curvesstack{\mathcal{C}\!\mathit{urves}}
\def\Polarizedstack{\mathcal{P}\!\mathit{olarized}}
\def\Complexesstack{\mathcal{C}\!\mathit{omplexes}}
% \Pic is the operator that assigns to X its picard group, usage \Pic(X)
% \Picardstack_{X/B} denotes the Picard stack of X over B
% \Picardfunctor_{X/B} denotes the Picard functor of X over B
\def\Pic{\mathop{\mathrm{Pic}}\nolimits}
\def\Picardstack{\mathcal{P}\!\mathit{ic}}
\def\Picardfunctor{\mathrm{Pic}}
\def\Deformationcategory{\mathcal{D}\!\mathit{ef}}

% --- appended 2026-10-12 (Chief): the project own definitions of the reference/slogan/history
% display environments, copied mechanically from the shared Stacks preamble lines 8-12
% (\usepackage{verbatim} plus the three \newenvironment lines, which the chapter preamble keeps
% verbatim from the project preamble). Nothing added, removed or redefined: these environments
% carry no mathematical content and the project itself typesets them as comments. Without these
% three lines a counted statement whose own text uses them fails with Environment slogan undefined.
\usepackage{verbatim}
\newenvironment{reference}{\comment}{\endcomment}
\newenvironment{slogan}{\comment}{\endcomment}
\newenvironment{history}{\comment}{\endcomment}
\usepackage{cleveref}
\crefname{originallemma}{original lemma}{original lemmas}
\Crefname{originallemma}{Original lemma}{Original lemmas}
\makeatletter
\newlabel{algebra-lemma-completion-flat}{{00MB}{source}{}{sourcepointer-1}{}}
\newlabel{algebra-lemma-completion-flat@cref}{{[originallemma][0][]00MB}{[1][0][]source}{}{}{}}
\makeatother
\crefname{originallemma}{original lemma}{original lemmas}
\Crefname{originallemma}{Original lemma}{Original lemmas}
\makeatletter
\newlabel{algebra-lemma-generic-flatness-Noetherian}{{051R}{source}{}{sourcepointer-2}{}}
\newlabel{algebra-lemma-generic-flatness-Noetherian@cref}{{[originallemma][0][]051R}{[1][0][]source}{}{}{}}
\makeatother
\crefname{originallemma}{original lemma}{original lemmas}
\Crefname{originallemma}{Original lemma}{Original lemmas}
\makeatletter
\newlabel{algebra-lemma-generic-flatness-reduced}{{051Z}{source}{}{sourcepointer-3}{}}
\newlabel{algebra-lemma-generic-flatness-reduced@cref}{{[originallemma][0][]051Z}{[1][0][]source}{}{}{}}
\makeatother
\crefname{originallemma}{original lemma}{original lemmas}
\Crefname{originallemma}{Original lemma}{Original lemmas}
\makeatletter
\newlabel{algebra-lemma-lift-projective}{{05CG}{source}{}{sourcepointer-4}{}}
\newlabel{algebra-lemma-lift-projective@cref}{{[originallemma][0][]05CG}{[1][0][]source}{}{}{}}
\makeatother
\crefname{originallemma}{original lemma}{original lemmas}
\Crefname{originallemma}{Original lemma}{Original lemmas}
\makeatletter
\newlabel{algebra-theorem-ffdescent-projectivity}{{05A9}{source}{}{sourcepointer-5}{}}
\newlabel{algebra-theorem-ffdescent-projectivity@cref}{{[originallemma][0][]05A9}{[1][0][]source}{}{}{}}
\makeatother
\crefname{originallemma}{original lemma}{original lemmas}
\Crefname{originallemma}{Original lemma}{Original lemmas}
\makeatletter
\newlabel{lemma-flat-pure-over-complete-projective}{{05DP}{source}{}{sourcepointer-6}{}}
\newlabel{lemma-flat-pure-over-complete-projective@cref}{{[originallemma][0][]05DP}{[1][0][]source}{}{}{}}
\makeatother
\title{A flat finite type algebra over a Noetherian ring with geometrically integral fibres is projective as a module}
\author{The Stacks Project Authors}
\date{stacks-project:flat.tex:local-snapshot-20261011; GNU FDL 1.2 or later}
\begin{document}
\maketitle
\noindent\textit{Editorial note.} All delivered mathematical text is the author's own, copied from the pinned Stacks Project chapter file flat.tex; no statement or proof step was rewritten, reordered or removed.\par\medskip
\noindent\textit{Editorial note.} Preamble provenance: the chapter's own retained minimal preamble (sources/stacks-flat/retained-preamble.tex, lines 1-85) is a verbatim excerpt of the native Xy-pic setup and of the native theorem declarations and mathematical macros, followed by the project's own reference/slogan/history display-environment definitions copied mechanically from the shared preamble; the wrapper class and the amsmath/amssymb/amsthm/graphicx base packages are supplied by the assembler, and extra\_packages is left empty because everything needed is already in the retained preamble. No macro was invented or redefined.\par\medskip
\noindent\textit{Editorial note.} Same-chapter pointers are located at their \textbackslash{}label lines inside the pinned chapter file; cross-chapter pointers are located in the project's own tag list staged in the snapshot (sources/stacks-chow/tags.txt), at the verbatim lines carrying their canonical tags.\par\medskip
\noindent\textit{Allowed standard prerequisites (editorial).}\begin{itemize}
\item Lifting projectivity along a nilpotent ideal. Algebra, label algebra-lemma-lift-projective, invoked at proof line 2202. Shows that if S/JS is projective over R/J with J the radical of I, then S/IS is projective over R/I because J/I is nilpotent and S/IS is flat, contradicting the maximality of I; hence I may be taken radical.
\item Generic flatness for a finite type algebra over a reduced Noetherian ring. Algebra, labels algebra-lemma-generic-flatness-Noetherian (proof line 2212) and algebra-lemma-generic-flatness-reduced (proof line 2214). Produces a nonzero f with S\_f free over R\_f, which enables the faithfully flat presentation of R by R\_f times its f-adic completion.
\item Faithfully flat descent of projectivity, flatness of a completion, and the criterion for a flat pure module over a complete local ring. Algebra, labels algebra-lemma-completion-flat (proof line 2222) and algebra-theorem-ffdescent-projectivity (proof line 2224); Same chapter flat.tex, label lemma-flat-pure-over-complete-projective (proof line 2227). Reduces projectivity of S over R to projectivity of the completed base change, and the criterion then applies because the fibre rings stay geometrically integral and the element f lies in the Jacobson radical of the complete ring.
\end{itemize}\medskip
\ifdefined\hypertarget\hypertarget{sourcepointer-1}{}\fi
\noindent\textit{Source reference.} Original lemma 00MB: sources/stacks-chow/tags.txt, source lines 817--817. Navigation pointer for \textbackslash{}ref\{algebra-lemma-completion-flat\} used inside the counted span (line 2222); the argument that consumes it is retained in the span.\par\medskip
\ifdefined\hypertarget\hypertarget{sourcepointer-2}{}\fi
\noindent\textit{Source reference.} Original lemma 051R: sources/stacks-chow/tags.txt, source lines 6544--6544. Navigation pointer for \textbackslash{}ref\{algebra-lemma-generic-flatness-Noetherian\} used inside the counted span (line 2212); the argument that consumes it is retained in the span.\par\medskip
\ifdefined\hypertarget\hypertarget{sourcepointer-3}{}\fi
\noindent\textit{Source reference.} Original lemma 051Z: sources/stacks-chow/tags.txt, source lines 6552--6552. Navigation pointer for \textbackslash{}ref\{algebra-lemma-generic-flatness-reduced\} used inside the counted span (line 2214); the argument that consumes it is retained in the span.\par\medskip
\ifdefined\hypertarget\hypertarget{sourcepointer-4}{}\fi
\noindent\textit{Source reference.} Original lemma 05CG: sources/stacks-chow/tags.txt, source lines 6919--6919. Navigation pointer for \textbackslash{}ref\{algebra-lemma-lift-projective\} used inside the counted span (line 2202); the argument that consumes it is retained in the span.\par\medskip
\ifdefined\hypertarget\hypertarget{sourcepointer-5}{}\fi
\noindent\textit{Source reference.} Original lemma 05A9: sources/stacks-chow/tags.txt, source lines 6842--6842. Navigation pointer for \textbackslash{}ref\{algebra-theorem-ffdescent-projectivity\} used inside the counted span (line 2224); the argument that consumes it is retained in the span.\par\medskip
\ifdefined\hypertarget\hypertarget{sourcepointer-6}{}\fi
\noindent\textit{Source reference.} Original lemma 05DP: sources/stacks-flat/pinned-original-source.tex, source lines 2145--2145. Navigation pointer for \textbackslash{}ref\{lemma-flat-pure-over-complete-projective\} used inside the counted span (line 2227); the argument that consumes it is retained in the span.\par\medskip
\begin{lemma}
\label{lemma-fibres-irreducible-flat-projective}
Let $R$ be a ring.
Let $R \to S$ be a ring map.
Assume
\begin{enumerate}
\item $R$ is Noetherian,
\item $R \to S$ is of finite type and flat, and
\item every fibre ring $S \otimes_R \kappa(\mathfrak p)$ is
geometrically integral over $\kappa(\mathfrak p)$.
\end{enumerate}
Then $S$ is projective as an $R$-module.
\end{lemma}

\begin{proof}
Consider the set
$$
\{I \subset R \mid S/IS\text{ not projective as }R/I\text{-module}\}
$$
We have to show this set is empty. To get a contradiction assume it is
nonempty. Then it contains a maximal element $I$.
Let $J = \sqrt{I}$ be its radical. If $I \not = J$, then
$S/JS$ is projective as a $R/J$-module, and $S/IS$ is flat over $R/I$
and $J/I$ is a nilpotent ideal in $R/I$. Applying
Algebra, Lemma \ref{algebra-lemma-lift-projective}
we see that $S/IS$ is a projective $R/I$-module, which is a contradiction.
Hence we may assume that $I$ is a radical ideal. In other words we
are reduced to proving the lemma in case $R$ is a reduced ring and
$S/IS$ is a projective $R/I$-module for every nonzero ideal $I$
of $R$.

\medskip\noindent
Assume $R$ is a reduced ring and $S/IS$ is a projective $R/I$-module
for every nonzero ideal $I$ of $R$. By generic flatness,
Algebra, Lemma \ref{algebra-lemma-generic-flatness-Noetherian}
(applied to a localization $R_g$ which is a domain) or the more general
Algebra, Lemma \ref{algebra-lemma-generic-flatness-reduced}
there exists a nonzero $f \in R$ such that $S_f$ is free as an
$R_f$-module. Denote $R^\wedge = \lim R/(f^n)$ the $(f)$-adic completion
of $R$. Note that the ring map
$$
R \longrightarrow R_f \times R^\wedge
$$
is a faithfully flat ring map, see
Algebra, Lemma \ref{algebra-lemma-completion-flat}.
Hence by faithfully flat descent of projectivity, see
Algebra, Theorem \ref{algebra-theorem-ffdescent-projectivity}
it suffices to prove that $S \otimes_R R^\wedge$ is a projective
$R^\wedge$-module. To see this we will use the criterion of
Lemma \ref{lemma-flat-pure-over-complete-projective}.
First of all, note that $S/fS = (S \otimes_R R^\wedge)/f(S \otimes_R R^\wedge)$
is a projective $R/(f)$-module and that $S \otimes_R R^\wedge$ is flat
and of finite type over $R^\wedge$ as a base change of such.
Next, suppose that $\mathfrak p^\wedge$ is a prime ideal
of $R^\wedge$. Let $\mathfrak p \subset R$ be the corresponding prime
of $R$. As $R \to S$ has geometrically integral fibre rings, the
same is true for the fibre rings of any base change. Hence
$\mathfrak q^\wedge = \mathfrak p^\wedge(S \otimes_R R^\wedge)$,
is a prime ideals lying over $\mathfrak p^\wedge$
and it is the unique associated prime of
$S \otimes_R \kappa(\mathfrak p^\wedge)$. Thus we win if
$f(S \otimes_R R^\wedge) + \mathfrak q^\wedge \not = S \otimes_R R^\wedge$.
This is true because $\mathfrak p^\wedge + fR^\wedge \not = R^\wedge$ as
$f$ lies in the Jacobson radical of the $f$-adically complete ring $R^\wedge$
and because $R^\wedge \to S \otimes_R R^\wedge$ is surjective on spectra
as its fibres are nonempty (irreducible spaces are nonempty).
\end{proof}

\end{document}
